\paragraph{作为\trMatchingCycle{}的\trVanishingCycle{}。}\label{sub:vanishing-as-matching}
取一个\trExact\trLefschetz\trBifibration，记号一如\eqref{eq:bifibration-diagram}。我们的头一个观察，是\trCriticalValue{}曲线 $\mathrm{Critv}(\varpi)\subset\scr{S}$ 的\trTopology{}及其到 $\scr{W}$ 上的\trProjection{}的一个简单推论。对每个 $w\in\scr{W}\setminus(\mathrm{Critv}(\Psi)\cup\mathrm{Fakev}(\Psi,\varpi))$，我们考察\trSurface{}$\scr{S}_w$ 与位于其中的、由点组成的有限构型 $\mathrm{Critv}(\varpi)\cap\scr{S}_w=\mathrm{Critv}(\varpi_w)$。由于该交是\trTransverse{}的，这一构型随 $w$ \trSmooth{}地变化。沿任一路径 $\beta\colon[0;1]\to\scr{W}\setminus(\mathrm{Critv}(\Psi)\cup\mathrm{Fakev}(\Psi,\varpi))$ 追随这一构型的运动，便可定义一个\textbf{\trBraid\trParallelTransport}（braid parallel transport）映射，它是一个保持\trOrientation{}的\trDiffeomorphism
\begin{equation*}
k_\beta\colon(\scr{S}_{\beta(0)},\mathrm{Critv}(\varpi_{\beta(0)}))
\longrightarrow(\scr{S}_{\beta(1)},\mathrm{Critv}(\varpi_{\beta(1)}))
\end{equation*}
在相差一个\trIsotopy{}的意义下唯一（固然可以把它取成辛的，但由于维数很低，由此并不会多得任何信息）。现在设 $\gamma\colon[0;1]\to\scr{W}$ 是 $\Psi$ 的一条\trVanishingPath，它同时避开 $\mathrm{Fakev}(\Psi,\varpi)$。\cref{thm:crit-local-embedding}告诉我们：当 $t\to1$ 时，$\mathrm{Critv}(\varpi)\cap\scr{S}_{\gamma(t)}$ 的诸点中有两点会在\trCovering{}$\psi|\mathrm{Critv}(\varpi)\colon\mathrm{Critv}(\varpi)\to\scr{W}$ 的一个\trBranchPoint{}处相互靠拢。取 $s$ 充分接近 $1$，在 $\scr{S}_{\gamma(s)}$ 中取连接这两点的短路径（在无\trOrientation{}的\trIsotopy{}下唯一）。它在 $k_{\gamma|[0;s]}$ 之下的\trPreimage{}是\trEmbedding{}路径 $\mu\colon[-1;1]\to\scr{S}_{\gamma(0)}$，满足 $\mu^{-1}(\mathrm{Critv}(\varpi_{\gamma(0)}))=\{-1;1\}$。我们称之为 $\gamma$ 相伴的\trMatchingPath，理由如下：
\begin{lemma}\label{thm:vanishing-as-matching-cycle}
$\mu$ 是 $\varpi_{\gamma(0)}$\footnote{译注：原书此处印作 $\varpi_{\mu(0)}$；疑应为 $\varpi_{\gamma(0)}$（$\varpi_w$ 仅对 $w\in\scr{W}$ 有定义，而 $\mu$ 是 $\scr{S}_{\gamma(0)}$ 中的路径；本证明首句亦以 $\varpi_{\gamma(s)}$ 称此\trObject）；此处径改。——译者注}的一条\trMatchingPath，并且对\trIsotopy{}$I$ 的适当选取，相伴的\trMatchingCycle{}$\Sigma_{\mu,I}$ 作为 $\scr{E}_{\gamma(0)}$ 中的一个带\trFrame{}的\trLagrange\trSphere，\trIsotopy{}于\trVanishingCycle{}$V_\gamma$。
\end{lemma}
\begin{proof}
只要有如下断言便已足够：在 $s$ 充分接近 $1$ 的 $\scr{E}_{\gamma(s)}$ 中，$\mathrm{Critv}(\varpi_{\gamma(s)})$ 的诸点里相互接近的那两点之间的短路径是一条\trMatchingPath，并且所得的\trMatchingCycle{}就是 $\gamma|[s;1]$ 的\trVanishingCycle。原陈述随即可经\trParallelTransport{}把一切从 $\scr{E}_{\gamma(s)}$ 送回 $\scr{E}_{\gamma(0)}$ 而得到。
考察落在 $\scr{E}_{\gamma(1)}$ 中的 $\Psi$ 的那个唯一\trCriticalPoint。我们在该点及其像附近引进\trHolomorphic{}局部坐标，使其具有\cref{thm:bifibration-normal-form}所述的性质。如\eqref{eq:symplectic-picard-lefschetz}的证明中或\cref{thm:matching-concatenation-twist}的证明中那样，我们可以局部修改 $\scr{E}$ 上的\trSymplecticForm{}使其在这些坐标中成为标准的。我们还可以\trDeformation{}路径 $\gamma$，使得在 $\Psi(x)\in\scr{W}$ 附近的局部坐标中，当 $t$ 充分接近 $1$ 时 $\gamma(t)=1-t$。考察短的\trVanishingPath{}$\gamma|[s;1]$，$s$ 充分接近 $1$。其相伴的\trVanishingCycle{}为
\begin{equation*}
V_{\gamma|[s;1]}=\sqrt{1-s}\,S^n\subset\CC^{n+1}
\end{equation*}
一如\cref{thm:toy-model-thimble}。另一方面，同一个球也是 $\scr{S}_{\gamma(s)}$ 中连接 $\varpi_{\gamma(s)}$ 的两个\trCriticalValue{}（$y_1=1-s$、$y_2=\pm\sqrt{1-s}$）的直线段所对应的\trMatchingCycle；验证这一点毫无困难，因为在这一情形中，\trMatchingCycle{}可以由朴素程序 \eqref{eq:naive-matching-cycle} 来定义。最后，如在\cref{thm:matching-concatenation-twist}的证明中那样，用\cref{thm:moser-deformation-embedding}便可回到原来的\trSymplecticForm。当然，对该形式而言，\trMatchingCycle{}与\trVanishingCycle{}的严格重合不再成立，但该论证仍为这两者提供了一个带\trFrame{}的\trIsotopy。
\end{proof}
\begin{remark}\label{thm:matching-path-isotopy-remark}
$V_\gamma$ 的\trIsotopy{}类只依赖于 $\gamma$ 在全体\trVanishingPath{}所成空间中的\trIsotopy{}类；特别地，若我们把 $\gamma$ 移过 $\mathrm{Fakev}(\Psi,\varpi)$ 的一点，它并不改变。然而，\trBraid\trParallelTransport{}通常会改变，从而相伴的\trMatchingPath{}$\mu$ 亦随之改变。这其间并无矛盾，因为不同的\trMatchingPath{}可以给出同一个\trMatchingCycle。事实上，\cref{thm:auroux-matching-example} 所提到的那一现象的诸实例，正是以这种方式来自最简单（\trGeneric）的\trFakeCriticalValue——在该处，曲线 $\mathrm{Critv}(\varpi)$\footnote{译注：原书此处印作 $\mathrm{Critv}(\pi)$（$\pi$ 上无杠）；疑应为 $\mathrm{Critv}(\varpi)$（\trFakeCriticalValue{}集 $\mathrm{Fakev}(\Psi,\varpi)$ 的定义所涉曲线即 $\mathrm{Critv}(\varpi)\subset\scr{S}$，而 $\pi$ 在\trBifibration{}记号中未定义）；此处径改。——译者注}有尖点或\trNode。
\end{remark}
\begin{remark}\label{thm:matching-cycle-history-remark}
如前文已经提到的，\trMatchingCycle{}系由\trDonaldson{}引入（未发表）。\trLagrange\trConnectedSum——特别在二维情形——在 [109] 中得到研究。较近以来，它在\trMirrorSymmetry{}中获得了一定的重要性，见 [140] 与 [139]；最后提到的那篇文章包含有\eqref{eq:matching-connected-sum}的一个版本。\cref{sub:vanishing-as-matching}的内容，对熟悉辫单值技术的读者——例如熟悉\trAuroux{}的\trBranchedCover{}[11] 与\trDonaldson{}的铅笔之间的关系的读者——将是显然的。
\end{remark}
